% GENERATED FILE. Edit canonical lessons, then run npm run generate:books.
% canonical-source-hash: fnv1a64:c83e5cb1d994d0c4
% canonical-lessons: GE-C213-eifersuechtig, GE-C213-erfolgreich, GE-C213-erkaeltet, GE-C213-fantastisch, GE-C213-feucht

\chapter{Jealous, Successful, Having a cold, Fantastic, Damp}
\label{ch:ge-jealous-successful-having-a-cold-fantastic-damp}
\hlchaptermodality{213}

\begin{chapteropening}
\textbf{By the end of this chapter:} \emph{I can say jealous, successful, having a cold, fantastic and damp.}

Complete the last lesson of chapter 213: I can say jealous, successful, having a cold, fantastic and damp.
\end{chapteropening}

\section[eifersüchtig]{eifersüchtig — jealous}

\label{lesson:GE-C213-eifersuechtig}



\begin{quote}
Before the new one: say the German for urgent, then the German for lonely.
\end{quote}

\subsection*{You'll want to know: eifersüchtig}
\textbf{eifersüchtig} — \textquotedblleft{}jealous\textquotedblright{}.

\begin{cousinweb}
\textbf{eifersüchtig} is the word for \textquotedblleft{}zeal\textquotedblright{} + the word for \textquotedblleft{}sickness, craving\textquotedblright{}, Old High German \emph{suht} (EtymDB 2.1), + \emph{-ig}: sick with zeal, jealous. English \emph{sick} is a cousin of that word.
\end{cousinweb}

\begin{grammarlens}[title={What you've built}]
One more word for this chapter.
\end{grammarlens}

\subsection*{Your turn}
\begin{itemize}
\raggedright
  \item \emph{Say it:} \emph{eifersüchtig}
  \item \emph{Say it:} \emph{eifersüchtig}, once more
  \item \emph{Say it:} say \emph{eifersüchtig}
  \item \emph{Recall:} say the German for urgent, then the German for lonely, then say \emph{eifersüchtig} again
\end{itemize}

\begin{tcolorbox}[breakable,colback=teal!4,colframe=teal!35!black,title={Before you move on}]
What does eifersüchtig mean? (\textquotedblleft{}Jealous\textquotedblright{}.) Say it once more.
\end{tcolorbox}
\section[erfolgreich]{erfolgreich — successful}

\label{lesson:GE-C213-erfolgreich}



\begin{quote}
Before the new one: say the German for lonely, then the German for jealous.
\end{quote}

\subsection*{You'll want to know: erfolgreich}
\textbf{erfolgreich} — \textquotedblleft{}successful\textquotedblright{}.

\begin{cousinweb}
\textbf{erfolgreich} is the word for \textquotedblleft{}success\textquotedblright{}, from a verb \textquotedblleft{}to follow, to result\textquotedblright{}, + \emph{reich}, \textquotedblleft{}rich\textquotedblright{} (EtymDB 2.1): rich in success. The \textquotedblleft{}follow\textquotedblright{} verb is the same as English \emph{follow}.
\end{cousinweb}

\begin{grammarlens}[title={What you've built}]
One more word for this chapter.
\end{grammarlens}

\subsection*{Your turn}
\begin{itemize}
\raggedright
  \item \emph{Say it:} \emph{erfolgreich}
  \item \emph{Say it:} \emph{erfolgreich}, once more
  \item \emph{Say it:} say \emph{erfolgreich}
  \item \emph{Recall:} say the German for lonely, then the German for jealous, then say \emph{erfolgreich} again
\end{itemize}

\begin{tcolorbox}[breakable,colback=teal!4,colframe=teal!35!black,title={Before you move on}]
What does erfolgreich mean? (\textquotedblleft{}Successful\textquotedblright{}.) Say it once more.
\end{tcolorbox}
\section[erkältet]{erkältet — having a cold}

\label{lesson:GE-C213-erkaeltet}



\begin{quote}
Before the new one: say the German for jealous, then the German for successful.
\end{quote}

\subsection*{You'll want to know: erkältet}
\textbf{erkältet} — \textquotedblleft{}having a cold\textquotedblright{}.

\begin{cousinweb}
\textbf{erkältet} is the past participle beside \emph{Erkältung}: \emph{er-} + the word for \textquotedblleft{}cold\textquotedblright{}, English \emph{cold}: chilled. (EtymDB 2.1 has Middle High German \emph{erkelten}.)
\end{cousinweb}

\begin{grammarlens}[title={What you've built}]
One more word for this chapter.
\end{grammarlens}

\subsection*{Your turn}
\begin{itemize}
\raggedright
  \item \emph{Say it:} \emph{erkältet}
  \item \emph{Say it:} \emph{erkältet}, once more
  \item \emph{Say it:} say \emph{erkältet}
  \item \emph{Recall:} say the German for jealous, then the German for successful, then say \emph{erkältet} again
\end{itemize}

\begin{tcolorbox}[breakable,colback=teal!4,colframe=teal!35!black,title={Before you move on}]
What does erkältet mean? (\textquotedblleft{}Having a cold\textquotedblright{}.) Say it once more.
\end{tcolorbox}
\section[fantastisch]{fantastisch — fantastic}

\label{lesson:GE-C213-fantastisch}



\begin{quote}
Before the new one: say the German for successful, then the German for having a cold.
\end{quote}

\subsection*{You'll want to know: fantastisch}
\textbf{fantastisch} — \textquotedblleft{}fantastic\textquotedblright{}.

\begin{cousinweb}
\textbf{fantastisch} $\leftarrow$ Latin \emph{phantasticus} $\leftarrow$ Greek \emph{phantastikós}, \textquotedblleft{}able to imagine, imaginary\textquotedblright{} (EtymDB 2.1) — borrowed.
\end{cousinweb}

\begin{grammarlens}[title={What you've built}]
One more word for this chapter.
\end{grammarlens}

\subsection*{Your turn}
\begin{itemize}
\raggedright
  \item \emph{Say it:} \emph{fantastisch}
  \item \emph{Say it:} \emph{fantastisch}, once more
  \item \emph{Say it:} say \emph{fantastisch}
  \item \emph{Recall:} say the German for successful, then the German for having a cold, then say \emph{fantastisch} again
\end{itemize}

\begin{tcolorbox}[breakable,colback=teal!4,colframe=teal!35!black,title={Before you move on}]
What does fantastisch mean? (\textquotedblleft{}Fantastic\textquotedblright{}.) Say it once more.
\end{tcolorbox}
\section[feucht]{feucht — damp}

\label{lesson:GE-C213-feucht}



\begin{quote}
Before the new one: say the German for having a cold, then the German for fantastic.
\end{quote}

\subsection*{You'll want to know: feucht}
\textbf{feucht} — \textquotedblleft{}damp\textquotedblright{}.

\begin{cousinweb}
\textbf{feucht} $\leftarrow$ Old High German \emph{fūht} $\leftarrow$ Proto-Germanic \emph{*funhtijaz}, \textquotedblleft{}humid, damp\textquotedblright{} — inherited. (EtymDB 2.1)
\end{cousinweb}

\begin{grammarlens}[title={What you've built}]
One more word for this chapter.
\end{grammarlens}

\subsection*{Your turn}
\begin{itemize}
\raggedright
  \item \emph{Say it:} \emph{feucht}
  \item \emph{Say it:} \emph{feucht}, once more
  \item \emph{Say it:} say \emph{feucht}
  \item \emph{Recall:} say the German for having a cold, then the German for fantastic, then say \emph{feucht} again
\end{itemize}

\begin{tcolorbox}[breakable,colback=teal!4,colframe=teal!35!black,title={Before you move on}]
What does feucht mean? (\textquotedblleft{}Damp\textquotedblright{}.) Say it once more.
\end{tcolorbox}
